% Options for packages loaded elsewhere
\PassOptionsToPackage{unicode}{hyperref}
\PassOptionsToPackage{hyphens}{url}
%
\documentclass[
]{article}
\usepackage{amsmath,amssymb}
\usepackage{lmodern}
\usepackage{iftex}
\ifPDFTeX
  \usepackage[T1]{fontenc}
  \usepackage[utf8]{inputenc}
  \usepackage{textcomp} % provide euro and other symbols
\else % if luatex or xetex
  \usepackage{unicode-math}
  \defaultfontfeatures{Scale=MatchLowercase}
  \defaultfontfeatures[\rmfamily]{Ligatures=TeX,Scale=1}
\fi
% Use upquote if available, for straight quotes in verbatim environments
\IfFileExists{upquote.sty}{\usepackage{upquote}}{}
\IfFileExists{microtype.sty}{% use microtype if available
  \usepackage[]{microtype}
  \UseMicrotypeSet[protrusion]{basicmath} % disable protrusion for tt fonts
}{}
\makeatletter
\@ifundefined{KOMAClassName}{% if non-KOMA class
  \IfFileExists{parskip.sty}{%
    \usepackage{parskip}
  }{% else
    \setlength{\parindent}{0pt}
    \setlength{\parskip}{6pt plus 2pt minus 1pt}}
}{% if KOMA class
  \KOMAoptions{parskip=half}}
\makeatother
\usepackage{xcolor}
\IfFileExists{xurl.sty}{\usepackage{xurl}}{} % add URL line breaks if available
\IfFileExists{bookmark.sty}{\usepackage{bookmark}}{\usepackage{hyperref}}
\hypersetup{
  pdftitle={Final C2O Strategy Audit Report},
  pdfauthor={C2O Coursework Submission},
  hidelinks,
  pdfcreator={LaTeX via pandoc}}
\urlstyle{same} % disable monospaced font for URLs
\usepackage{longtable,booktabs,array}
\usepackage{calc} % for calculating minipage widths
% Correct order of tables after \paragraph or \subparagraph
\usepackage{etoolbox}
\makeatletter
\patchcmd\longtable{\par}{\if@noskipsec\mbox{}\fi\par}{}{}
\makeatother
% Allow footnotes in longtable head/foot
\IfFileExists{footnotehyper.sty}{\usepackage{footnotehyper}}{\usepackage{footnote}}
\makesavenoteenv{longtable}
\setlength{\emergencystretch}{3em} % prevent overfull lines
\providecommand{\tightlist}{%
  \setlength{\itemsep}{0pt}\setlength{\parskip}{0pt}}
\setcounter{secnumdepth}{-\maxdimen} % remove section numbering
\ifLuaTeX
  \usepackage{selnolig}  % disable illegal ligatures
\fi

\title{Final C2O Strategy Audit Report}
\author{C2O Coursework Submission}
\date{}

\begin{document}
\maketitle

\hypertarget{final-c2o-strategy-audit-report}{%
\section{Final C2O Strategy Audit
Report}\label{final-c2o-strategy-audit-report}}

\textbf{Coursework strategy:} close-to-open equity strategy using the
promoted volatility-scaled overnight return target.\\
\textbf{Final target:}
\texttt{volatility\_scaled\_overnight\_return\ =\ overnight\_next\ /\ (vol20\ /\ sqrt(252))}.\\
\textbf{Volatility denominator:} trailing 20-day close-to-close
volatility shifted by one trading day.\\
\textbf{Development cutoff:} 2024-12-31. 2025-2026 is not used for
development and is treated as held out for marker evaluation.\\
\textbf{Headline 250M result:} 3.93\% net annual return, 4.90\% net
volatility, 0.811 net Sharpe, and -10.19\% maximum drawdown.\\
\textbf{Primary sources:} \texttt{outputs/performance\_summary.csv},
\texttt{outputs/report\_ready/final\_strategy\_config.md},
\texttt{outputs/report\_ready/final\_strategy\_selection\_memo.md}, and
\texttt{outputs/report\_ready/final\_reproduction\_log.md}.

\hypertarget{1-executive-summary}{%
\subsection{1. Executive Summary}\label{1-executive-summary}}

The final promoted C2O strategy is not the original baseline. The
original baseline remains the starting point and is archived in
\texttt{outputs/baseline\_archive/}. The archive readme records an
original baseline pipeline Sharpe of about 0.374 at 250M and a logged
experiment-harness baseline Sharpe of about 0.379 at 250M. The
difference is disclosed in
\texttt{outputs/baseline\_archive/baseline\_readme.md} and is treated as
an implementation/logging distinction rather than a change in economic
thesis.

The final strategy uses the same point-in-time feature set as the
baseline, but changes the supervised learning target from raw overnight
return to a volatility-scaled overnight return. Economically, this asks
the model to rank stocks by risk-adjusted overnight opportunity rather
than by raw overnight magnitude. That is a narrower and more defensible
improvement than adding new signals after seeing results: the feature
set, universe rules, costs, borrow treatment, participation cap, and
development cutoff remain frozen.

At 250M AUM the final strategy earns 3.93\% net annual return with
4.90\% net volatility, a net Sharpe of 0.811, and a maximum drawdown of
-10.19\%. These numbers come from
\texttt{outputs/performance\_summary.csv}. The result is modest rather
than spectacular. It is credible precisely because it is reported net of
commission, auction slippage, and borrow costs; because the 5\% ADV20
participation cap is enforced in the positions files; and because the
final audit acknowledges that performance degrades under larger capital.

The strongest caution is capacity. At 50M the strategy uses nearly the
full target gross exposure and records a net Sharpe of 0.850. At 250M
average gross exposure falls to 70.43\% because the participation cap
binds on almost every trading day. At 1B it falls to 22.82\%, and the
net Sharpe falls to 0.266. This is not presented as a scalable
institutional anomaly at arbitrary size. It is a small-to-mid-capacity
overnight implementation whose economics are materially shaped by
trading constraints.

The second major caution is costs. The 250M gross Sharpe is 2.427, but
commission, auction slippage, and borrow reduce it to 0.811. Auction
slippage is the largest single degradation item. The report therefore
avoids gross-return marketing and treats implementation as part of the
strategy, not an afterthought.

\hypertarget{final-performance-by-aum}{%
\subsubsection{Final Performance by
AUM}\label{final-performance-by-aum}}

\begin{longtable}[]{@{}lllrlrrllr@{}}
\toprule
AUM & Net annual return & Net vol & Net Sharpe & Max drawdown & Gross
Sharpe & Avg turnover & Avg gross exposure & IC mean & IC t-stat \\
\midrule
\endhead
\$50m & 4.30\% & 5.11\% & 0.85 & -15.50\% & 3.071 & 1.998 & 99.90\% &
2.82\% & 9.465 \\
\$250m & 3.93\% & 4.90\% & 0.811 & -10.19\% & 2.427 & 1.409 & 70.43\% &
2.82\% & 9.465 \\
\$1bn & 0.67\% & 2.64\% & 0.266 & -16.06\% & 1.239 & 0.456 & 22.82\% &
2.82\% & 9.465 \\
\bottomrule
\end{longtable}

Source: \texttt{outputs/performance\_summary.csv}. Basket selection
rule, participation cap, cost schedule, and target version are
documented in \texttt{outputs/report\_ready/final\_strategy\_config.md}.

\hypertarget{baseline-versus-final-candidate-at-250m}{%
\subsubsection{Baseline Versus Final Candidate at
250M}\label{baseline-versus-final-candidate-at-250m}}

\begin{longtable}[]{@{}lllrlrrrrrllrl@{}}
\toprule
Strategy & Net annual return & Net vol & Net Sharpe & Max DD & Gross
Sharpe & Commission drag & Slippage drag & Borrow drag & Turnover & Avg
gross exp. & IC mean & IC t-stat & Worst 12m \\
\midrule
\endhead
Original baseline (experiment harness) & 1.69\% & 4.71\% & 0.379 &
-24.58\% & 1.729 & 0.309 & 0.928 & 0.112 & 1.158 & 57.89\% & 2.73\% &
8.363 & -12.98\% \\
Final volatility-scaled target & 3.93\% & 4.90\% & 0.811 & -10.19\% &
2.427 & 0.361 & 1.084 & 0.17 & 1.409 & 70.43\% & 2.82\% & 9.465 &
-10.09\% \\
\bottomrule
\end{longtable}

Source: \texttt{outputs/report\_ready/final\_candidate\_comparison.csv};
baseline archive explanation in
\texttt{outputs/baseline\_archive/baseline\_readme.md}.

\hypertarget{2-data-return-construction-and-point-in-time-controls}{%
\subsection{2. Data, Return Construction, and Point-in-Time
Controls}\label{2-data-return-construction-and-point-in-time-controls}}

The strategy is built around overnight returns from close on day
\texttt{t} to open on day \texttt{t+1}. The final modelling target is
the next overnight return scaled by trailing realized volatility. The
denominator is \texttt{vol20\ /\ sqrt(252)}, where \texttt{vol20} is a
20-trading-day close-to-close volatility estimate shifted by one trading
day. This shift is the critical point-in-time control: the target
transformation cannot see the close-to-close return that is still
unknown at the 15:50 ET decision time on day \texttt{t}.

The point-in-time audit in
\texttt{outputs/report\_ready/volatility\_scaled\_target\_point\_in\_time\_audit.md}
is the basis for promoting the candidate. It verifies that the
volatility denominator uses trailing information, is shifted, does not
use future overnight returns, does not use future close/high/low
observations, does not use full-sample volatility, and does not
introduce 2025+ data. The audit also confirms that the target
transformation changes only the training target and not the live feature
set.

The return construction was reconciled against the available price
fields. The frozen reconciliation file reports the following checks:

\begin{longtable}[]{@{}ll@{}}
\toprule
Metric & Value \\
\midrule
\endhead
tolerance & 1.000e-08 \\
stock\_days\_checked & 5,469,968 \\
fail\_count & 0 \\
fail\_fraction & 0 \\
median\_abs\_residual & 0.000e+00 \\
p99\_abs\_residual & 2.220e-16 \\
max\_abs\_residual & 4.441e-16 \\
\bottomrule
\end{longtable}

Source: \texttt{outputs/return\_reconciliation\_summary.csv}.

\hypertarget{earnings-timing-controls}{%
\subsubsection{Earnings Timing
Controls}\label{earnings-timing-controls}}

Earnings data are a classic look-ahead risk in close-to-open strategies.
The frozen audit examples include one after-market-close case, one
before-market-open case, and a same-day example. The rules preserve
decision-time availability: before-market events can be known before
that day's close, while after-market events are not available before
that close.

\begin{longtable}[]{@{}lrllllll@{}}
\toprule
Ticker & Instrument & Reporting date & Timing & Raw event date &
Strategy trading date & Excluded decision dates & Rule applied \\
\midrule
\endhead
MOS & 18 & 2010-01-05 & after & 2010-01-05 & 2010-01-05 & 2010-01-04,
2010-01-05, 2010-01-06 & After-market event is not known before the
close; use provided strategy trading date. \\
RPM & 1275 & 2010-01-06 & before & 2010-01-06 & 2010-01-05 & 2010-01-04,
2010-01-05, 2010-01-06 & Before-market event is known before that day's
close; use provided strategy trading date. \\
MOS & 18 & 2010-01-05 & after & 2010-01-05 & 2010-01-05 & 2010-01-04,
2010-01-05, 2010-01-06 & After-market event is not known before the
close; use provided strategy trading date. \\
\bottomrule
\end{longtable}

Source: \texttt{outputs/earnings\_timing\_examples.csv} and
\texttt{outputs/report\_ready/earnings\_timing\_audit.md}.

The hand-checkable examples matter because an earnings feature can
accidentally become a future-information feature if all reporting dates
are treated identically. The final audit does not rely on a generic
statement that the data are lagged; it exports concrete examples with
\texttt{raw\_event\_date}, \texttt{strat\_trading\_date}, excluded
decision dates, the rule applied, and a text explanation.

\hypertarget{short-interest-timing-controls}{%
\subsubsection{Short-Interest Timing
Controls}\label{short-interest-timing-controls}}

Short-interest data are lagged relative to economic reality and can be
dangerous if the model is allowed to use a value before it could have
been published or vendor-processed. The submitted pipeline treats the
provided \texttt{all\_data} short-interest proxies as already
point-in-time transformed according to the coursework data package, and
then applies an additional one-trading-day strategy lag before using
\texttt{dsi}, \texttt{dtcn}, and \texttt{ddtcn} in features or
borrow-tier assignment. The audit examples are:

\begin{longtable}[]{@{}lrllrrrl@{}}
\toprule
Ticker & Instrument & Source feature date & Decision date using lag1 &
DSI & DTCN & DDTCN & Rule applied \\
\midrule
\endhead
HLT & 1 & 2015-01-27 & 2015-01-28 & 0.004074 & 1.66949 & -0.245679 &
Short-interest proxies are read from the daily all\_data panel, then
shifted one trading day in add\_point\_in\_time\_features. \\
HLT & 1 & 2015-01-28 & 2015-01-29 & 0.004074 & 1.66949 & -0.245679 &
Short-interest proxies are read from the daily all\_data panel, then
shifted one trading day in add\_point\_in\_time\_features. \\
HLT & 1 & 2015-01-29 & 2015-01-30 & 0.004074 & 1.66949 & -0.245679 &
Short-interest proxies are read from the daily all\_data panel, then
shifted one trading day in add\_point\_in\_time\_features. \\
\bottomrule
\end{longtable}

Source: \texttt{outputs/short\_interest\_lag\_examples.csv} and
\texttt{outputs/report\_ready/short\_interest\_lag\_audit.md}.

The limitation is explicit: the frozen outputs do not contain raw FINRA
settlement and publication dates. The report therefore does not claim
independent reconstruction of the original short-interest publication
lag. It claims only what is supported by the repository evidence: the
strategy verifies decision-time use of the provided point-in-time series
and adds a one-trading-day decision lag. This is also listed in
\texttt{outputs/report\_ready/report\_missing\_items.md}.

\hypertarget{feature-inventory-and-observability}{%
\subsubsection{Feature Inventory and
Observability}\label{feature-inventory-and-observability}}

The final feature set is unchanged from the baseline. The report-ready
feature inventory records the formula, data source, required lag, and
whether the feature is observable at 15:50 ET on day \texttt{t}. A
compact extract is shown below; the complete table is in
\texttt{outputs/report\_ready/feature\_inventory.csv}.

\begin{longtable}[]{@{}lllll@{}}
\toprule
Feature & Formula & Data source & Required lag & Observable at 15:50 ET
day t \\
\midrule
\endhead
rank\_r\_on\_1d & r\_overnight = adj\_open\_t / adj\_close\_\{t-1\} - 1
& prices.parquet: adjusted open, adjusted close & 0 trading days;
Open\_t is known after the 09:30 auction & Yes \\
rank\_r\_on\_5d & rolling mean of r\_overnight over 5 sessions,
min\_periods=3 & prices.parquet: adjusted open, adjusted close & 0
trading days; uses overnight returns available by 15:50 ET & Yes \\
rank\_r\_id\_1d & r\_intraday\_\{t-1\} = adj\_close\_\{t-1\} /
adj\_open\_\{t-1\} - 1 & prices.parquet: adjusted open, adjusted close &
1 trading day & Yes \\
rank\_ret\_cc\_5d & adj\_close\_\{t-1\} / adj\_close\_\{t-6\} - 1 &
prices.parquet: adjusted close & 1 trading day & Yes \\
rank\_ret\_cc\_20d & adj\_close\_\{t-1\} / adj\_close\_\{t-21\} - 1 &
prices.parquet: adjusted close & 1 trading day & Yes \\
rank\_ret\_cc\_60d & adj\_close\_\{t-1\} / adj\_close\_\{t-61\} - 1 &
prices.parquet: adjusted close & 1 trading day & Yes \\
rank\_vol20 & annualised rolling std of close-to-close returns over 20
sessions, shifted 1 day & prices.parquet: adjusted close & 1 trading day
& Yes \\
rank\_vol60 & annualised rolling std of close-to-close returns over 60
sessions, shifted 1 day & prices.parquet: adjusted close & 1 trading day
& Yes \\
rank\_amihud20 & rolling mean over 20 sessions of abs(r\_close\_close) /
dollar\_volume, shifted 1 day & prices.parquet: adjusted close, volume &
1 trading day & Yes \\
rank\_adv20\_log & log(1 + trailing 20-day average dollar volume),
shifted 1 day & prices.parquet: adjusted close, volume & 1 trading day &
Yes \\
rank\_market\_cap\_log & log(1 + market\_cap\_\{t-1\}) & prices.parquet:
market\_cap & 1 trading day & Yes \\
rank\_piot\_norm\_lag1 & piot\_norm\_\{t-1\} & all\_data.parquet & 1
trading day & Yes \\
rank\_gross\_profit\_margin\_lag1 & gross\_profit\_margin\_\{t-1\} &
all\_data.parquet & 1 trading day & Yes \\
rank\_asset\_turnover\_ratio\_lag1 & asset\_turnover\_ratio\_\{t-1\} &
all\_data.parquet & 1 trading day & Yes \\
rank\_net\_debt\_to\_equity\_lag1 & net\_debt\_to\_equity\_\{t-1\} &
all\_data.parquet & 1 trading day & Yes \\
rank\_price\_to\_book\_lag1 & price\_to\_book\_\{t-1\} &
all\_data.parquet & 1 trading day & Yes \\
rank\_ev\_to\_ebit\_lag1 & ev\_to\_ebit\_\{t-1\} & all\_data.parquet & 1
trading day & Yes \\
rank\_sue\_lag1 & sue\_\{t-1\} & all\_data.parquet & 1 trading day &
Yes \\
rank\_deps\_lag1 & deps\_\{t-1\} & all\_data.parquet & 1 trading day &
Yes \\
rank\_reps1\_lag1 & reps1\_\{t-1\} & all\_data.parquet & 1 trading day &
Yes \\
rank\_repsf4\_lag1 & repsf4\_\{t-1\} & all\_data.parquet & 1 trading day
& Yes \\
rank\_final\_score\_lag1 & final\_score\_\{t-1\} &
cheapness\_scores.parquet & 1 trading day & Yes \\
rank\_score\_velocity\_lag1 & score\_velocity\_\{t-1\} &
cheapness\_scores.parquet & 1 trading day & Yes \\
rank\_momentum\_score\_lag1 & momentum\_score\_\{t-1\} &
cheapness\_scores.parquet & 1 trading day & Yes \\
rank\_dsi\_lag1 & dsi\_\{t-1\}; daily short-interest proxy after
publication/vendor lag already in source & all\_data.parquet & 1 trading
day after source availability & Yes, if source obeys Section 2.1.3 lag
convention \\
rank\_dtcn\_lag1 & dtcn\_\{t-1\}; daily days-to-cover proxy after
publication/vendor lag already in source & all\_data.parquet & 1 trading
day after source availability & Yes, if source obeys Section 2.1.3 lag
convention \\
rank\_ddtcn\_lag1 & ddtcn\_\{t-1\}; lagged change/stress proxy after
publication/vendor lag already in source & all\_data.parquet & 1 trading
day after source availability & Yes, if source obeys Section 2.1.3 lag
convention \\
rank\_industry\_return\_lag1 & industry\_return\_\{t-1\} &
all\_data.parquet & 1 trading day & Yes \\
\bottomrule
\end{longtable}

Source: \texttt{outputs/report\_ready/feature\_inventory.csv}.

\hypertarget{3-universe-construction-and-capacity-filters}{%
\subsection{3. Universe Construction and Capacity
Filters}\label{3-universe-construction-and-capacity-filters}}

The final universe, eligibility filters, basket rule, and participation
cap are frozen in
\texttt{outputs/report\_ready/final\_strategy\_config.md}. The final
report treats these implementation choices as part of the strategy
definition, because the overnight signal is only investable if it can be
traded near the close and open with realistic size.

The final capacity constraint is a 5\% ADV20 participation cap.
Positions are clipped against the cap before final weights and dollar
positions are recorded. The frozen positions files include
\texttt{adv20\_dollar}, \texttt{participation\_rate},
\texttt{cap\_dollar}, and \texttt{cap\_binding\_flag}, which allows the
reported exposure reduction to be checked from position-level evidence
rather than from summary claims alone.

\hypertarget{universe-evolution}{%
\subsubsection{Universe Evolution}\label{universe-evolution}}

\begin{longtable}[]{@{}rrllr@{}}
\toprule
Year & Universe Count & Ref Date & Median Year Start Mcap & Mid Year
Exits \\
\midrule
\endhead
2010 & 986 & 2009-12-31 & \$2.46bn & 0 \\
2011 & 986 & 2010-12-31 & \$3.25bn & 0 \\
2012 & 988 & 2011-12-30 & \$3.25bn & 0 \\
2013 & 988 & 2012-12-31 & \$3.86bn & 0 \\
2014 & 987 & 2013-12-31 & \$5.42bn & 0 \\
2015 & 986 & 2014-12-31 & \$6.07bn & 0 \\
2016 & 987 & 2015-12-31 & \$6.06bn & 0 \\
2017 & 988 & 2016-12-30 & \$7.07bn & 0 \\
2018 & 988 & 2017-12-29 & \$8.73bn & 0 \\
2019 & 988 & 2018-12-31 & \$7.47bn & 0 \\
2020 & 989 & 2019-12-31 & \$9.62bn & 0 \\
2021 & 990 & 2020-12-31 & \$11.62bn & 0 \\
2022 & 988 & 2021-12-31 & \$14.16bn & 0 \\
2023 & 988 & 2022-12-30 & \$11.50bn & 0 \\
2024 & 989 & 2023-12-29 & \$13.06bn & 0 \\
\bottomrule
\end{longtable}

Source: \texttt{outputs/universe\_summary.csv}.

\hypertarget{eligibility-exclusions}{%
\subsubsection{Eligibility Exclusions}\label{eligibility-exclusions}}

\begin{longtable}[]{@{}ll@{}}
\toprule
Eligibility reason & Stock-days \\
\midrule
\endhead
OK & 3,194,145 \\
MCAP\_FAIL & 1,596,321 \\
ADV\_FAIL & 513,267 \\
EARN\_WINDOW & 116,925 \\
PRICE\_FAIL & 26,307 \\
VOL\_FAIL & 22,458 \\
DATA\_FAIL & 1,000 \\
\bottomrule
\end{longtable}

Source: \texttt{outputs/eligibility\_summary.csv}.

\hypertarget{capacity-evidence}{%
\subsubsection{Capacity Evidence}\label{capacity-evidence}}

\begin{longtable}[]{@{}llrllllll@{}}
\toprule
AUM & Avg gross exposure used & Cap-binding days & Fraction cap-binding
days & Avg per-stock abs position & Max per-stock abs position & Avg
participation & Max participation & Position-days at cap \\
\midrule
\endhead
\$50m & 99.90\% & 3464 & 91.79\% & \$1.00m & \$7.40m & 1.99\% & 5.00\% &
16.16\% \\
\$250m & 70.43\% & 3773 & 99.97\% & \$3.53m & \$98.01m & 3.76\% & 5.00\%
& 61.68\% \\
\$1bn & 22.82\% & 3773 & 99.97\% & \$4.58m & \$414.68m & 3.96\% & 5.00\%
& 67.87\% \\
\bottomrule
\end{longtable}

Source: \texttt{outputs/report\_ready/position\_capacity\_summary.csv}.

The capacity table is one of the most important controls in the report.
At 250M, the strategy is not simply choosing to run lower risk. It is
capacity constrained: cap-binding days are nearly the whole sample,
average gross exposure used is 70.43\%, and 61.68\% of position-days sit
at the cap. At 1B, average gross exposure drops to 22.82\% and 67.87\%
of position-days are capped. That explains why the 1B result is much
weaker even though the underlying ranking signal remains the same.

The frozen outputs report fixed auction slippage and realized
participation, but do not export a separate square-root market-impact
estimate at the 5\% cap. The report therefore does not introduce a new
impact model. This limitation is listed in
\texttt{outputs/report\_ready/report\_missing\_items.md}.

\hypertarget{4-borrow-proxy-and-short-leg-treatment}{%
\subsection{4. Borrow Proxy and Short-Leg
Treatment}\label{4-borrow-proxy-and-short-leg-treatment}}

The short leg is not treated as free. The final strategy applies tiered
borrow costs based on the hard-to-borrow proxy documented in
\texttt{outputs/report\_ready/borrow\_proxy\_methodology.md}. The proxy
uses short-interest and borrow-stress style variables from the processed
data. The final implementation applies borrow cost tiers to short
positions; it does not hard-exclude short candidates during selection.

The hard-to-borrow proxy is deterministic and uses lagged short-interest
stress fields from the point-in-time panel. Tier A is the default borrow
tier at 40 bps per annum. Tier B applies when
\texttt{dsi\_lag1\ \textgreater{}=\ 0.08}, or
\texttt{dtcn\_lag1\ \textgreater{}=\ 5.0}, or
\texttt{ddtcn\_lag1\ \textgreater{}=\ 1.0}, and is charged at 200 bps
per annum. Tier C applies when
\texttt{dsi\_lag1\ \textgreater{}=\ 0.15}, or
\texttt{dtcn\_lag1\ \textgreater{}=\ 10.0}, or both
\texttt{dsi\_lag1\ \textgreater{}=\ 0.10} and
\texttt{ddtcn\_lag1\ \textgreater{}=\ 1.5}, and is charged at 800 bps
per annum. Borrow is charged daily on short notional as
\texttt{borrow\_rate\_annual\ /\ 252}. Source:
\texttt{outputs/report\_ready/borrow\_proxy\_methodology.md}.

\begin{longtable}[]{@{}llrl@{}}
\toprule
Borrow tier & Proxy threshold & Annual borrow rate & Final treatment \\
\midrule
\endhead
A & Default tier when Tier B/C thresholds are not met & 40 bps & Cost
applied to short notional \\
B & \texttt{dsi\_lag1\ \textgreater{}=\ 0.08}, or
\texttt{dtcn\_lag1\ \textgreater{}=\ 5.0}, or
\texttt{ddtcn\_lag1\ \textgreater{}=\ 1.0} & 200 bps & Cost applied to
short notional \\
C & \texttt{dsi\_lag1\ \textgreater{}=\ 0.15}, or
\texttt{dtcn\_lag1\ \textgreater{}=\ 10.0}, or both
\texttt{dsi\_lag1\ \textgreater{}=\ 0.10} and
\texttt{ddtcn\_lag1\ \textgreater{}=\ 1.5} & 800 bps & Cost applied to
short notional \\
\bottomrule
\end{longtable}

The final audit reports that 94,035 raw short position-days were
observed, of which 57,039 are Tier B or Tier C. This is 60.66\% of raw
short position-days. The same file reports that the fraction of raw
short signal affected by borrow treatment is 0.00\% because borrow
treatment is cost-only in the final configuration, not a selection
filter. Source:
\texttt{outputs/report\_ready/short\_signal\_affected\_by\_borrow.csv}.

\hypertarget{borrow-cost-degradation}{%
\subsubsection{Borrow Cost Degradation}\label{borrow-cost-degradation}}

\begin{longtable}[]{@{}lrrrrrr@{}}
\toprule
AUM & Gross Sharpe & Commission drag & Auction slippage drag & Borrow
drag & Net Sharpe & Avg daily borrow bps \\
\midrule
\endhead
\$50m & 3.071 & 0.493 & 1.48 & 0.249 & 0.85 & 0.503 \\
\$250m & 2.427 & 0.361 & 1.084 & 0.17 & 0.811 & 0.331 \\
\$1bn & 1.239 & 0.218 & 0.654 & 0.1 & 0.266 & 0.105 \\
\bottomrule
\end{longtable}

Source: \texttt{outputs/report\_ready/borrow\_sharpe\_degradation.csv}.

\hypertarget{borrow-tier-summary}{%
\subsubsection{Borrow Tier Summary}\label{borrow-tier-summary}}

\begin{longtable}[]{@{}llrlllr@{}}
\toprule
AUM & Tier & Short position-days & Share & Avg short notional & Total
borrow cost & Avg daily borrow bps \\
\midrule
\endhead
\$50m & A & 36996 & 39.34\% & \$1.10m & \$645.08k & 0.503 \\
\$50m & B & 36226 & 38.52\% & \$948.06k & \$2.73m & 0.503 \\
\$50m & C & 20813 & 22.13\% & \$926.10k & \$6.12m & 0.503 \\
\$250m & A & 36996 & 39.34\% & \$4.02m & \$2.36m & 0.331 \\
\$250m & B & 36226 & 38.52\% & \$3.40m & \$9.78m & 0.331 \\
\$250m & C & 20813 & 22.13\% & \$2.89m & \$19.10m & 0.331 \\
\$1bn & A & 36996 & 39.34\% & \$5.48m & \$3.22m & 0.105 \\
\$1bn & B & 36226 & 38.52\% & \$4.16m & \$11.97m & 0.105 \\
\$1bn & C & 20813 & 22.13\% & \$3.70m & \$24.44m & 0.105 \\
\bottomrule
\end{longtable}

Source: \texttt{outputs/report\_ready/borrow\_tier\_summary.csv}.

The borrow table is deliberately reported in position-days, not just
dollars, because a strategy can hide borrow risk by netting costs into
aggregate returns. At 250M, Tier A accounts for 39.34\% of short
position-days, Tier B for 38.52\%, and Tier C for 22.13\%. Tier C has
fewer position-days than the other tiers but a larger annualized cost
assumption, so it accounts for a material share of total borrow cost.

The external validation limitation is also explicit. No independent
locate-fee or borrow-fee file is present in the frozen outputs, so the
report cannot claim external validation of the proxy thresholds. It
reports the proxy definition, thresholds, tier shares, and cost
consequences instead. This is a sceptical rather than promotional
treatment of the short leg.

\hypertarget{5-alpha-model-and-target-engineering}{%
\subsection{5. Alpha Model and Target
Engineering}\label{5-alpha-model-and-target-engineering}}

The final model is a cross-sectional overnight alpha model. Its key
design choice is the promoted target:
\texttt{overnight\_next\ /\ (vol20\ /\ sqrt(252))}. This target is
economically interpretable. Raw overnight returns can be dominated by
volatile stocks; volatility scaling rewards names whose expected
overnight opportunity is large relative to their recent close-to-close
risk. The change is therefore a target-engineering decision, not a new
data source.

The final feature set remains unchanged from the baseline. This matters
for look-ahead control. If the Sharpe improvement had come from adding
new features, the audit burden would be wider. Here the main question is
narrower: whether the denominator in the transformed target is
point-in-time valid. The dedicated audit passed that question.

The final 250M IC mean is 2.82\% with an IC t-stat of 9.465. The
baseline experiment-harness IC mean is 2.73\% with an IC t-stat of
8.363. Source:
\texttt{outputs/report\_ready/final\_candidate\_comparison.csv}. The
improvement is not only a realized volatility artifact; the ranking
statistic is slightly stronger as well. However, the report does not
overstate this. The realized performance improvement is larger than the
IC improvement, indicating that risk scaling and realized portfolio
volatility also contribute.

\hypertarget{feature-ablation-at-250m}{%
\subsubsection{Feature Ablation at
250M}\label{feature-ablation-at-250m}}

\begin{longtable}[]{@{}lllrllrlrl@{}}
\toprule
Variant & Removed group & IC mean & IC t-stat & Net annual return & Net
vol & Net Sharpe & Max DD & Avg turnover & Avg gross exposure \\
\midrule
\endhead
full\_model & none & 2.82\% & 9.465 & 3.93\% & 4.90\% & 0.811 & -10.19\%
& 1.409 & 70.43\% \\
remove\_return/reversal & return/reversal & 0.79\% & 2.106 & -4.80\% &
5.02\% & -0.954 & -55.66\% & 1.011 & 50.53\% \\
remove\_risk/liquidity/size & risk/liquidity/size & 2.85\% & 11.364 &
4.49\% & 4.82\% & 0.934 & -8.48\% & 1.605 & 80.27\% \\
remove\_fundamental/value/quality & fundamental/value/quality & 2.92\% &
9.784 & 4.14\% & 4.75\% & 0.878 & -10.38\% & 1.382 & 69.10\% \\
remove\_earnings/revision & earnings/revision & 2.75\% & 9.177 & 4.05\%
& 4.85\% & 0.841 & -10.19\% & 1.407 & 70.34\% \\
remove\_short-interest/borrow-stress & short-interest/borrow-stress &
2.85\% & 9.52 & 4.99\% & 5.22\% & 0.959 & -11.32\% & 1.474 & 73.69\% \\
remove\_industry-return & industry-return & 2.82\% & 9.469 & 3.84\% &
4.90\% & 0.795 & -10.19\% & 1.408 & 70.40\% \\
\bottomrule
\end{longtable}

Source: \texttt{outputs/feature\_ablation\_summary.csv} and
\texttt{outputs/report\_ready/feature\_ablation\_audit.md}.

The ablation audit is mixed, which is useful. Removing return/reversal
features sharply damages performance: the 250M net Sharpe falls from
0.811 to -0.954. That is the clearest evidence that the core alpha is an
overnight return/reversal mechanism. Removing risk/liquidity/size and
short-interest/borrow-stress improves 250M Sharpe in the frozen ablation
table, so the report should not claim that every feature group adds
standalone alpha. This does not automatically make those groups useless:
some variables can act as risk, capacity, crowding, turnover, or
cost-control information rather than pure alpha predictors. The right
conclusion is narrower: return/reversal is the dominant alpha source,
while the remaining groups have mixed incremental value under the frozen
model.

The final strategy was not retuned during ablation. The ablation is
therefore an audit of feature dependence under the frozen modelling
setup, not a fresh search over hyperparameters.

\hypertarget{6-portfolio-construction-and-cost-model}{%
\subsection{6. Portfolio Construction and Cost
Model}\label{6-portfolio-construction-and-cost-model}}

The final strategy selects a daily long and short basket according to
the frozen basket selection rule in
\texttt{outputs/report\_ready/final\_strategy\_config.md}. The positions
are held overnight from the decision close to the next open. The final
positions files are exported for 50M, 250M, and 1B AUM and include the
fields needed to audit score, target weight, final clipped weight,
dollar position, cap binding, gross PnL, commission, slippage, borrow
cost, net PnL, and borrow tier.

The cost model includes commission, auction slippage, and borrow. This
report treats auction slippage as an essential part of the close/open
implementation, not as a sensitivity to be added later. At 250M, gross
Sharpe is 2.427, but commission drag is 0.361, auction slippage drag is
1.084, and borrow drag is 0.170, leaving net Sharpe of 0.811. Source:
\texttt{outputs/report\_ready/borrow\_sharpe\_degradation.csv}.

The same cost table shows why gross results are misleading. At 50M gross
Sharpe is 3.071 and net Sharpe is 0.850; at 1B gross Sharpe is 1.239 and
net Sharpe is 0.266. Costs are not a cosmetic adjustment; they decide
whether the strategy is economically interesting.

\hypertarget{7-empirical-results}{%
\subsection{7. Empirical Results}\label{7-empirical-results}}

At the chosen reporting size of 250M, the final strategy has a net
annual return of 3.93\%, net volatility of 4.90\%, net Sharpe of 0.811,
and max drawdown of -10.19\%. These are the headline numbers in
\texttt{outputs/performance\_summary.csv} and
\texttt{outputs/report\_ready/final\_reproduction\_log.md}.

The result is credible but not grand. A 3.93\% annual return is not
enough to survive careless implementation or inflated fees. It becomes
interesting because volatility is controlled, because the drawdown is
materially smaller than the archived baseline comparison, and because
the result survives cost and capacity constraints at 250M. It is still
not a claim of live profitability in 2025-2026. Those years are held out
for marker evaluation.

\hypertarget{year-by-year-final-results-at-250m}{%
\subsubsection{Year-by-Year Final Results at
250M}\label{year-by-year-final-results-at-250m}}

\begin{longtable}[]{@{}rllrlll@{}}
\toprule
Year & Annual return & Annual vol & Annual Sharpe & Max DD & IC mean &
Avg gross exposure \\
\midrule
\endhead
2010 & -10.09\% & 3.42\% & -3.094 & -10.19\% & -3.48\% & 42.46\% \\
2011 & 21.39\% & 5.63\% & 3.474 & -2.08\% & 6.22\% & 67.54\% \\
2012 & 6.70\% & 2.80\% & 2.333 & -1.68\% & 4.67\% & 56.93\% \\
2013 & 5.61\% & 2.38\% & 2.307 & -1.17\% & 7.22\% & 64.02\% \\
2014 & 2.57\% & 3.66\% & 0.712 & -5.01\% & 4.23\% & 80.35\% \\
2015 & 1.58\% & 2.31\% & 0.691 & -1.73\% & 3.75\% & 57.44\% \\
2016 & 1.13\% & 1.97\% & 0.578 & -1.83\% & 4.21\% & 38.38\% \\
2017 & -0.46\% & 1.38\% & -0.331 & -1.60\% & 2.79\% & 42.84\% \\
2018 & 2.82\% & 3.01\% & 0.939 & -1.99\% & 2.17\% & 69.11\% \\
2019 & 5.85\% & 3.53\% & 1.63 & -1.38\% & 2.98\% & 83.78\% \\
2020 & -3.15\% & 9.23\% & -0.3 & -9.08\% & 0.81\% & 86.51\% \\
2021 & 2.70\% & 5.47\% & 0.515 & -5.70\% & -2.11\% & 77.36\% \\
2022 & 28.21\% & 8.10\% & 3.111 & -2.86\% & 6.29\% & 95.46\% \\
2023 & 3.99\% & 5.95\% & 0.687 & -5.88\% & 2.82\% & 95.42\% \\
2024 & -4.13\% & 6.17\% & -0.652 & -8.30\% & -0.29\% & 98.81\% \\
\bottomrule
\end{longtable}

Source:
\texttt{outputs/report\_ready/baseline\_vs\_candidate\_yearly.csv}.

The final strategy has positive annual returns in 11 out of 15 calendar
years in the frozen 250M yearly table. Its weakest year is 2010, with
annual return -10.09\%. Its strongest year is 2022, with annual return
28.21\%. These figures come from
\texttt{outputs/report\_ready/baseline\_vs\_candidate\_yearly.csv}.

\hypertarget{quantstats-tear-sheet}{%
\subsubsection{QuantStats Tear Sheet}\label{quantstats-tear-sheet}}

A QuantStats HTML tear sheet for the final 250M strategy is saved at
\texttt{outputs/quantstats\_250m.html}. Figure caption for any use of
that tear sheet: AUM = 250M; basket selection rule = frozen final basket
rule in \texttt{outputs/report\_ready/final\_strategy\_config.md};
participation cap = 5\% ADV20; cost schedule = commission, auction
slippage, and tiered borrow costs from the frozen final configuration;
target version = volatility-scaled overnight return.

\hypertarget{8-robustness-feature-ablation-and-stress-windows}{%
\subsection{8. Robustness, Feature Ablation, and Stress
Windows}\label{8-robustness-feature-ablation-and-stress-windows}}

The stress-window analysis compares the archived logged baseline and the
final volatility-scaled target in three difficult windows: late 2018,
2020 Q1, and the 2022 drawdown. The final result is not uniformly better
on every metric, which is important. In late 2018 the final strategy has
a positive net return but a slightly worse max drawdown than the
baseline. In 2020 Q1 and 2022, it is materially stronger in the frozen
outputs.

\hypertarget{stress-windows-at-250m}{%
\subsubsection{Stress Windows at 250M}\label{stress-windows-at-250m}}

\begin{longtable}[]{@{}lllrllrr@{}}
\toprule
Window & Strategy & Net return & Annualized Sharpe & Max DD & Avg gross
exposure & Avg turnover & Days \\
\midrule
\endhead
covid\_q1\_2020 & Baseline & -2.45\% & -0.743 & -6.77\% & 54.99\% & 1.1
& 62 \\
covid\_q1\_2020 & Final & 2.24\% & 1.02 & -3.05\% & 87.97\% & 1.759 &
62 \\
drawdown\_2022 & Baseline & -7.16\% & -0.788 & -13.36\% & 93.64\% &
1.873 & 251 \\
drawdown\_2022 & Final & 28.09\% & 3.111 & -2.86\% & 95.46\% & 1.909 &
251 \\
late\_2018 & Baseline & -0.95\% & -2.093 & -1.08\% & 30.89\% & 0.618 &
63 \\
late\_2018 & Final & 0.17\% & 0.235 & -1.30\% & 67.95\% & 1.359 & 63 \\
\bottomrule
\end{longtable}

Source:
\texttt{outputs/report\_ready/baseline\_vs\_candidate\_stress\_windows.csv}.

The stress table supports promotion, but not triumphalism. The 2022
window is particularly strong for the final strategy, with 28.09\% net
return and 3.111 annualized Sharpe. That helps the full-sample result. A
sceptical marker should ask whether that is a single lucky sub-period.
The year-by-year table partially addresses this: the final strategy is
positive in 11 of 15 years, but not in every year, and 2010 remains a
materially negative year.

The feature ablation table provides a different robustness check. The
strategy depends heavily on return/reversal information. It is less
clearly helped by every other group. This supports a focused economic
interpretation rather than a broad claim that all available data sources
are equally useful.

\hypertarget{9-limitations-and-held-out-evaluation-risk}{%
\subsection{9. Limitations and Held-Out Evaluation
Risk}\label{9-limitations-and-held-out-evaluation-risk}}

The main limitations are capacity, cost realism, borrow validation, and
future-period uncertainty.

Capacity is the cleanest limitation because it is measured directly. At
1B, the strategy cannot deploy enough capital under the 5\% ADV20 cap.
Average gross exposure used is only 22.82\%, and the net Sharpe is
0.266. The negative intuition is not that the alpha disappears
completely; it is that the tradable expression becomes too constrained.
A rank signal with insufficient executable notional is not a scalable
portfolio.

Costs are also central. The final 250M gross Sharpe of 2.427 would look
impressive if shown alone. Net Sharpe is 0.811 after commission, auction
slippage, and borrow. The difference is the implementation reality of
trading a close-to-open basket.

Borrow honesty is incomplete but disclosed. The final output includes
tiered borrow costs and short-position tier shares. It does not include
external borrow-fee validation, and the strategy does not change
selection based on borrow tier. That means hard-to-borrow exposure
affects net returns through cost, not through avoided trades.

Short-interest timing is controlled at the processed-panel level, but
the repository does not include raw publication-date evidence. The final
report therefore avoids claiming more than the files support.

Finally, 2025-2026 is held out for marker evaluation. The report does
not claim that the final 2010-2024 result necessarily persists. The
correct interpretation is narrower: under the frozen data cutoff and
final implementation, the volatility-scaled target is a point-in-time
valid improvement over the archived baseline.

\hypertarget{10-section-7-sceptical-marker-questions}{%
\subsection{10. Section 7 Sceptical-Marker
Questions}\label{10-section-7-sceptical-marker-questions}}

\textbf{Look-ahead.} The promoted target passed a point-in-time audit.
\texttt{vol20} is trailing and shifted by one trading day; the feature
set is unchanged; and the final reproduction log confirms the cached
final panel maxes out at 2024-12-31. Earnings and short-interest audits
provide hand-checkable examples. The remaining short-interest limitation
is source-panel publication timing, disclosed in
\texttt{outputs/report\_ready/report\_missing\_items.md}.

\textbf{Statistical robustness.} The 250M IC mean is 2.82\% with t-stat
9.465. The final strategy is positive in 11 of 15 calendar years.
Feature ablation shows the model is materially dependent on
return/reversal features; it is not robust to removing that group. The
stress-window results are favorable in 2020 Q1 and 2022 but less clean
in late 2018.

\textbf{Capacity and execution honesty.} The report shows 50M, 250M, and
1B AUM. The 5\% ADV20 cap is enforced and exported at position level.
The 1B result is weak because exposure cannot be deployed under the cap.
Auction slippage is the largest Sharpe drag at 250M and is reported
separately.

\textbf{Borrow honesty.} The report includes tiered borrow costs,
borrow-tier position-day shares, and gross-to-net Sharpe degradation. It
also states that external borrow-fee validation is not present and that
borrow treatment is cost-only rather than a hard exclusion.

\textbf{Reporting integrity.} The original baseline is archived, the
final strategy is frozen, final outputs are generated under
\texttt{outputs/} and \texttt{outputs/report\_ready/}, and
reproducibility is documented.
\texttt{outputs/report\_ready/final\_reproduction\_log.md} records that
\texttt{make\ reproduce} and \texttt{make\ test} passed, with 13 tests
passing.

\hypertarget{11-conclusion}{%
\subsection{11. Conclusion}\label{11-conclusion}}

The final strategy should be presented as a disciplined improvement over
the baseline, not as a discovered magic formula. The promoted target is
economically sensible and point-in-time valid: it ranks overnight
opportunities by expected return per unit of recent realized volatility.
At 250M, the final net Sharpe of 0.811 is materially higher than the
archived baseline comparison and the max drawdown is materially smaller.

The strategy is still capacity limited. It is most credible around the
50M to 250M scale and becomes weak at 1B under the 5\% ADV20 cap. It is
also highly sensitive to execution costs, especially auction slippage.
These caveats are not side notes; they are part of the final conclusion.

The recommended final submission is therefore the frozen
volatility-scaled overnight return strategy, with 2025-2026 left
untouched for external evaluation.

\hypertarget{appendix-a-frozen-final-configuration-crosswalk}{%
\subsection{Appendix A. Frozen Final Configuration
Crosswalk}\label{appendix-a-frozen-final-configuration-crosswalk}}

The following configuration table is copied from the frozen strategy
configuration rather than reconstructed from code. Its purpose is to
make the submitted report auditable without asking the marker to infer
the strategy settings from multiple files.

\begin{longtable}[]{@{}lll@{}}
\toprule
Configuration item & Frozen setting & Source \\
\midrule
\endhead
Final strategy name &
\texttt{E\_target\_transform\ /\ volatility\_scaled\_overnight\_return}
& \texttt{outputs/report\_ready/final\_strategy\_config.md} \\
Target & \texttt{volatility\_scaled\_overnight\_return} &
\texttt{outputs/report\_ready/final\_strategy\_config.md} \\
Target definition & \texttt{overnight\_next\ /\ (vol20\ /\ sqrt(252))} &
\texttt{outputs/report\_ready/final\_strategy\_config.md} \\
Overnight return definition &
\texttt{adj\_open\_(t+1)\ /\ adj\_close\_t\ -\ 1} &
\texttt{outputs/report\_ready/final\_strategy\_config.md} \\
Volatility denominator & trailing 20-day close-to-close volatility,
annualised, shifted by one trading day &
\texttt{outputs/report\_ready/final\_strategy\_config.md} \\
Feature set & unchanged from the original baseline point-in-time feature
set & \texttt{outputs/report\_ready/final\_strategy\_config.md} \\
Universe & annual top-1000 US common-equity universe by market
capitalisation &
\texttt{outputs/report\_ready/final\_strategy\_config.md} \\
Universe freeze & each year uses the last trading day of the previous
year & \texttt{outputs/report\_ready/final\_strategy\_config.md} \\
Minimum history & 252 price observations &
\texttt{outputs/report\_ready/final\_strategy\_config.md} \\
Minimum price & \$5.00 &
\texttt{outputs/report\_ready/final\_strategy\_config.md} \\
Minimum ADV20 & \$10,000,000 &
\texttt{outputs/report\_ready/final\_strategy\_config.md} \\
Volatility band & 5\% to 120\% annualised &
\texttt{outputs/report\_ready/final\_strategy\_config.md} \\
Earnings exclusion & +/- 1 trading day &
\texttt{outputs/report\_ready/final\_strategy\_config.md} \\
Long basket & top 3\% of eligible names by alpha score &
\texttt{outputs/report\_ready/final\_strategy\_config.md} \\
Short basket & bottom 3\% of eligible names by alpha score &
\texttt{outputs/report\_ready/final\_strategy\_config.md} \\
Minimum basket size & 15 names per side &
\texttt{outputs/report\_ready/final\_strategy\_config.md} \\
Weighting & equal weighting within each side, subject to capacity caps &
\texttt{outputs/report\_ready/final\_strategy\_config.md} \\
Dollar neutrality & long and short books matched after capacity sizing &
\texttt{outputs/report\_ready/final\_strategy\_config.md} \\
Participation cap & 5.0\% of ADV20 per name &
\texttt{outputs/report\_ready/final\_strategy\_config.md} \\
AUM levels & 50M, 250M, 1B &
\texttt{outputs/report\_ready/final\_strategy\_config.md} \\
Commission & 0.5 bps per leg, 1.0 bps round trip &
\texttt{outputs/report\_ready/final\_strategy\_config.md} \\
Auction slippage & 1.5 bps per leg, 3.0 bps round trip &
\texttt{outputs/report\_ready/final\_strategy\_config.md} \\
Non-borrow round-trip cost & 4.0 bps on notional turnover &
\texttt{outputs/report\_ready/final\_strategy\_config.md} \\
Borrow Tier A & 40 bps per annum &
\texttt{outputs/report\_ready/final\_strategy\_config.md} \\
Borrow Tier B & 200 bps per annum &
\texttt{outputs/report\_ready/final\_strategy\_config.md} \\
Borrow Tier C & 800 bps per annum &
\texttt{outputs/report\_ready/final\_strategy\_config.md} \\
Borrow application & charged daily on short notional as annual rate /
252 & \texttt{outputs/report\_ready/final\_strategy\_config.md} \\
Hard borrow exclusion & none in final strategy &
\texttt{outputs/report\_ready/final\_strategy\_config.md} \\
Start date & 2010-01-01 &
\texttt{outputs/report\_ready/final\_strategy\_config.md} \\
Development cutoff & 2024-12-31 &
\texttt{outputs/report\_ready/final\_strategy\_config.md} \\
Random seed & 42; deterministic implementation &
\texttt{outputs/report\_ready/final\_strategy\_config.md} \\
\bottomrule
\end{longtable}

This table is also the caption disclosure standard for any final
strategy figure. Where a figure is strategy-specific, the caption must
identify AUM, basket rule, 5\% ADV20 participation cap,
commission/slippage/borrow cost schedule, and volatility-scaled target
version.

\hypertarget{appendix-b-coursework-brief-crosswalk}{%
\subsection{Appendix B. Coursework Brief
Crosswalk}\label{appendix-b-coursework-brief-crosswalk}}

This crosswalk states where the report answers the coursework brief. It
is included to reduce ambiguity for marking: each answer is tied to a
frozen output file, not a later calculation.

\begin{longtable}[]{@{}lll@{}}
\toprule
Brief area & Report answer & Frozen source files \\
\midrule
\endhead
Section 2: return construction & Close-to-open return uses adjusted
close at day \texttt{t} and adjusted open at day \texttt{t+1};
reconciliation reports zero failures under the exported tolerance. &
\texttt{outputs/return\_reconciliation\_summary.csv} \\
Section 2: point-in-time controls & Volatility denominator is trailing
and shifted; target transformation changes only the label, not features.
&
\texttt{outputs/report\_ready/volatility\_scaled\_target\_point\_in\_time\_audit.md},
\texttt{outputs/report\_ready/final\_strategy\_config.md} \\
Section 2: earnings timing & AMC, BMO, and same-day timing examples are
exported with event date, strategy trading date, excluded decision
dates, rule, and explanation. &
\texttt{outputs/earnings\_timing\_examples.csv},
\texttt{outputs/report\_ready/earnings\_timing\_audit.md} \\
Section 2: short-interest timing & Processed short-interest proxies are
shifted one trading day before strategy use; raw publication-date
limitation is disclosed. &
\texttt{outputs/short\_interest\_lag\_examples.csv},
\texttt{outputs/report\_ready/short\_interest\_lag\_audit.md},
\texttt{outputs/report\_ready/report\_missing\_items.md} \\
Section 3: universe construction & Annual top-1000 universe is frozen
from the prior year-end reference date; universe counts and median
market cap are exported by year. &
\texttt{outputs/universe\_summary.csv},
\texttt{outputs/report\_ready/final\_strategy\_config.md} \\
Section 3: capacity filters & Minimum price, ADV20, volatility band,
earnings window, and required data filters are stated; exclusion
stock-days are exported by reason. &
\texttt{outputs/eligibility\_summary.csv},
\texttt{outputs/report\_ready/final\_strategy\_config.md} \\
Section 3: implementation capacity & 5\% ADV20 cap is enforced and
audited through average exposure, cap-binding days, position size, and
participation summaries. &
\texttt{outputs/report\_ready/position\_capacity\_summary.csv},
\texttt{outputs/positions\_50m.parquet},
\texttt{outputs/positions\_250m.parquet},
\texttt{outputs/positions\_1b.parquet} \\
Section 4: borrow proxy & Borrow treatment is tiered cost rather than
hard exclusion; tier rates, tier shares, and cost impact are reported. &
\texttt{outputs/report\_ready/borrow\_proxy\_methodology.md},
\texttt{outputs/report\_ready/borrow\_tier\_summary.csv},
\texttt{outputs/report\_ready/borrow\_sharpe\_degradation.csv} \\
Section 5: alpha target & Final target is risk-adjusted overnight
return, not raw overnight return. &
\texttt{outputs/report\_ready/final\_strategy\_config.md},
\texttt{outputs/report\_ready/final\_strategy\_selection\_memo.md} \\
Section 5: feature groups & Final feature set is unchanged; feature
formulas and observability are exported. &
\texttt{outputs/report\_ready/feature\_inventory.csv},
\texttt{outputs/report\_ready/feature\_inventory.md} \\
Section 5: IC evidence & IC mean and t-stat are reported for baseline
and final candidate, and by calendar year in the yearly table. &
\texttt{outputs/report\_ready/final\_candidate\_comparison.csv},
\texttt{outputs/report\_ready/baseline\_vs\_candidate\_yearly.csv} \\
Section 5: ablation & Final-strategy feature-group ablation is reported
at 50M, 250M, and 1B, with the report emphasizing the 250M table. &
\texttt{outputs/feature\_ablation\_summary.csv},
\texttt{outputs/report\_ready/feature\_ablation\_audit.md} \\
Section 6: portfolio construction & Long/short top-bottom 3\% baskets,
minimum 15 names per side, equal side weighting, dollar neutrality, and
cap clipping are documented. &
\texttt{outputs/report\_ready/final\_strategy\_config.md} \\
Section 6: cost model & Commission, auction slippage, and borrow costs
are included; gross-to-net Sharpe degradation is decomposed. &
\texttt{outputs/report\_ready/borrow\_sharpe\_degradation.csv},
\texttt{outputs/performance\_summary.csv} \\
Section 6: empirical results & Final 50M, 250M, and 1B performance is
reported net of costs, with yearly and stress-window checks. &
\texttt{outputs/performance\_summary.csv},
\texttt{outputs/report\_ready/baseline\_vs\_candidate\_yearly.csv},
\texttt{outputs/report\_ready/baseline\_vs\_candidate\_stress\_windows.csv} \\
Section 7: look-ahead scepticism & Report answers target, feature,
earnings, short-interest, and data-cutoff look-ahead risks separately. &
\texttt{outputs/report\_ready/volatility\_scaled\_target\_point\_in\_time\_audit.md},
\texttt{outputs/report\_ready/final\_reproduction\_log.md} \\
Section 7: statistical robustness & Report uses IC, yearly results,
ablation, and stress windows rather than only headline Sharpe. &
\texttt{outputs/report\_ready/final\_candidate\_comparison.csv},
\texttt{outputs/feature\_ablation\_summary.csv},
\texttt{outputs/report\_ready/baseline\_vs\_candidate\_stress\_windows.csv} \\
Section 7: execution honesty & Capacity cap and cost drag are presented
as central limitations. &
\texttt{outputs/report\_ready/position\_capacity\_summary.csv},
\texttt{outputs/report\_ready/borrow\_sharpe\_degradation.csv} \\
Section 7: borrow honesty & Borrow tier costs and the absence of
external validation are both disclosed. &
\texttt{outputs/report\_ready/borrow\_tier\_summary.csv},
\texttt{outputs/report\_ready/report\_missing\_items.md} \\
Section 7: reporting integrity & Baseline archived; final strategy
frozen; reproduction and tests logged. &
\texttt{outputs/baseline\_archive/baseline\_readme.md},
\texttt{outputs/report\_ready/final\_reproduction\_log.md} \\
\bottomrule
\end{longtable}

\hypertarget{appendix-c-frozen-output-manifest}{%
\subsection{Appendix C. Frozen Output
Manifest}\label{appendix-c-frozen-output-manifest}}

The final report deliberately uses a narrow source set. The table below
records how each source is used, so that the reported numbers can be
traced without rerunning experiments.

\begin{longtable}[]{@{}ll@{}}
\toprule
Frozen artifact & Use in report \\
\midrule
\endhead
\texttt{outputs/performance\_summary.csv} & final headline performance
by AUM, including net return, volatility, Sharpe, drawdown, turnover,
exposure, and IC \\
\texttt{outputs/report\_ready/final\_strategy\_config.md} & final
strategy definition, target, universe, eligibility filters, basket rule,
costs, borrow, cutoff, and seed \\
\texttt{outputs/report\_ready/final\_strategy\_selection\_memo.md} &
promotion rationale and comparison logic for the final candidate \\
\texttt{outputs/report\_ready/final\_strategy\_summary.md} &
final-strategy summary language and headline context \\
\texttt{outputs/baseline\_archive/baseline\_readme.md} & original
baseline preservation and baseline Sharpe disclosure \\
\texttt{outputs/report\_ready/final\_candidate\_comparison.csv} &
baseline versus final candidate comparison across AUM levels \\
\texttt{outputs/report\_ready/baseline\_vs\_candidate\_yearly.csv} &
calendar-year performance and IC stability \\
\texttt{outputs/report\_ready/baseline\_vs\_candidate\_stress\_windows.csv}
& late 2018, 2020 Q1, and 2022 stress-window comparison \\
\texttt{outputs/report\_ready/position\_capacity\_summary.csv} & average
gross exposure, cap-binding days, per-name positions, participation, and
capped position-days \\
\texttt{outputs/report\_ready/borrow\_sharpe\_degradation.csv} &
gross-to-net Sharpe degradation by commission, slippage, and borrow \\
\texttt{outputs/report\_ready/borrow\_tier\_summary.csv} & short
position-day counts, shares, notional, and borrow cost by tier \\
\texttt{outputs/report\_ready/short\_signal\_affected\_by\_borrow.csv} &
distinction between borrow cost impact and selection impact \\
\texttt{outputs/feature\_ablation\_summary.csv} & final-strategy
feature-group ablation results \\
\texttt{outputs/report\_ready/feature\_ablation\_audit.md} & explanation
of ablation method and no-retuning control \\
\texttt{outputs/report\_ready/feature\_inventory.csv} & formula, data
source, lag, and observability for features \\
\texttt{outputs/earnings\_timing\_examples.csv} & hand-checkable AMC,
BMO, and same-day earnings examples \\
\texttt{outputs/short\_interest\_lag\_examples.csv} & hand-checkable
short-interest lag examples \\
\texttt{outputs/return\_reconciliation\_summary.csv} & adjusted price
return reconciliation check \\
\texttt{outputs/universe\_summary.csv} & annual universe size and
year-start market-cap reference evidence \\
\texttt{outputs/eligibility\_summary.csv} & exclusion stock-days by
eligibility reason \\
\texttt{outputs/report\_ready/final\_reproduction\_log.md} & Python
version, commands, test result, final metrics, cutoff confirmation,
baseline archive confirmation \\
\texttt{outputs/quantstats\_250m.html} & final 250M tear sheet
artifact \\
\texttt{outputs/report\_ready/report\_missing\_items.md} & evidence
limitations that are not filled with unsupported estimates \\
\bottomrule
\end{longtable}

The manifest also sets a boundary on the report. If a claim is not
supported by one of these frozen files, it is either omitted or
disclosed as unavailable. This is why the report does not include
external borrow-fee validation or a new market-impact model at the
participation cap.

\hypertarget{file-traceability}{%
\subsection{File Traceability}\label{file-traceability}}

\begin{itemize}
\tightlist
\item
  Final strategy configuration:
  \texttt{outputs/report\_ready/final\_strategy\_config.md}
\item
  Final headline performance: \texttt{outputs/performance\_summary.csv}
\item
  Baseline archive note:
  \texttt{outputs/baseline\_archive/baseline\_readme.md}
\item
  Final comparison:
  \texttt{outputs/report\_ready/final\_candidate\_comparison.csv}
\item
  Yearly comparison:
  \texttt{outputs/report\_ready/baseline\_vs\_candidate\_yearly.csv}
\item
  Stress windows:
  \texttt{outputs/report\_ready/baseline\_vs\_candidate\_stress\_windows.csv}
\item
  Capacity summary:
  \texttt{outputs/report\_ready/position\_capacity\_summary.csv}
\item
  Borrow degradation:
  \texttt{outputs/report\_ready/borrow\_sharpe\_degradation.csv}
\item
  Borrow tiers: \texttt{outputs/report\_ready/borrow\_tier\_summary.csv}
\item
  Feature ablation: \texttt{outputs/feature\_ablation\_summary.csv}
\item
  Timing audits: \texttt{outputs/earnings\_timing\_examples.csv},
  \texttt{outputs/short\_interest\_lag\_examples.csv}
\item
  Reproducibility:
  \texttt{outputs/report\_ready/final\_reproduction\_log.md}
\item
  Missing evidence list:
  \texttt{outputs/report\_ready/report\_missing\_items.md}
\end{itemize}

\end{document}
